\documentclass[oneside,a4paper,14pt]{extarticle} %размер шрифта 14
\usepackage[T1,T2A]{fontenc}
\usepackage[
	a4paper,
	letterpaper,
	top=2cm,
	bottom=2cm,
	left=2.5cm,
	right=1.5cm
]{geometry}
\usepackage[utf8]{inputenc} %кодировка текста
\usepackage[russian]{babel} %поддержка русского языка
\usepackage{textcomp} %текстовые символы
\usepackage{indentfirst} %корректировка отступов
\usepackage{listings} %для кода !!!
\usepackage{xcolor}
\usepackage{float}
\usepackage{enumitem} %для нумерации буквами
\usepackage{graphicx} %работа с изображениям
\usepackage{mwe} % for blindtext and example-image-a in example
\usepackage{wrapfig}
\usepackage{caption}
\usepackage{amsmath}  % для формул и символов
\usepackage{amsfonts}
\usepackage{amsthm}
\usepackage{graphicx}
\usepackage[all]{xy}
\usepackage[breaklinks]{hyperref}

% размер кегля у заголоков разделов
\usepackage{titlesec}
\titleformat{\section}
{\normalsize\bfseries}
{\thesection} {1em}{}
\titleformat{\subsection}
{\normalsize\bfseries}
{\thesubsection} {1em}{}
\titleformat{\subsubsection}
{\normalsize\bfseries}
{\thesubsection} {1em}{}

\renewcommand\baselinestretch{1.33}\normalsize % межстрочный интервал
\setlength{\parindent}{1.25cm}
\usepackage{indentfirst}

% \setminted{fontsize=\small}

\begin{document}
	\newpage\thispagestyle{empty}
	\begin{center}
		МИНИСТЕРСТВО НАУКИ И ВЫСШЕГО ОБРАЗОВАНИЯ\\
			РОССИЙСКОЙ ФЕДЕРАЦИИ
			ФЕДЕРАЛЬНОЕ ГОСУДАРСТВЕННОЕ БЮДЖЕТНОЕ\\
			ОБРАЗОВАТЕЛЬНОЕ
			УЧРЕЖДЕНИЕ ВЫСШЕГО ОБРАЗОВАНИЯ\\
			«ВЯТСКИЙ ГОСУДАРСТВЕННЫЙ УНИВЕРСИТЕТ»\\
			Институт математики и информационных систем\\
			Факультет автоматики и вычислительной техники\\
			Кафедра электронных вычислительных машин
	\end{center}
	\vspace{20mm}

	\begin{center}
		Отчёт по лабораторной работе №2\\
		по дисциплине\\
		<<Промпт-инжиниринг>>\\
		<<Создание и редактирование визуального контента \\с помощью
промпт-инжиниринга>>\\
		<<Вариант 4>>\\
	\end{center}
	\vspace{48mm}

	\noindent
	Выполнил студент гр. ИВТб-1303-06-00 \hfill \rule[-0,5mm]{30mm}{0.15mm}\,/Гортоломей И.К./
	\newline
	Проверил доцент кафедры ЭВМ \hfill \rule[-0,5mm]{30mm}{0.15mm}\,/Колупаев А.В./

	\vfill
	\begin{center}
		Киров\\
		2026
	\end{center}

	\newpage
	\thispagestyle{empty}

	\section*{Цель работы:}
Освоение инженерных подходов к управлению генерацией и обработкой визуального
контента с использованием нейросетевых моделей. Основная цель - научиться
контролировать визуальный вывод моделей, обеспечивать стилистическую
согласованность (consistency), генерировать валидные синтетические данные для
обучения ML-алгоритмов, а также применять мультимодальные модели для извлечения
структурированных данных из неструктурированных визуальных источников.

	\section*{Используемые модели}
	\noindent
	\indent (1) gpt-image-1.5-high-fidelity \\
	\indent (2) gemini-3-pro-image-preview-2k (nano-banana-pro) \\
	\newline
	\noindent Данные модели предоставляются сайтом arena.ai, поэтому управление параметрами, такими как Seed, Steps, CFG Scale и др., на прямую невозможно.

	\setcounter{section}{1}

	\newpage
	\section*{Задание \thesection : Генерация изображения по текстовому запросу}
	\subsection{исходные данные}
	\indent \textbf{Центральный объект:} Умный светофор с камерами \\
	\indent \textbf{Окружение / Фон:} Перекресток мегаполиса под дождем \\
	\indent \textbf{Стиль изображения:} Кинематографичный кадр (как в кино) \\
	\indent \textbf{Освещение:} Ночное уличное, отражения на мокром асфальте \\
	\indent \textbf{Обязательные детали:} Объективы камер наблюдения, рамки выделения объектов (bounding boxes) вокруг машин
	\subsection{базовый промпт}
	\begin{verbatim}
					Умный светофор на перекрестке мегаполиса, ночь, дождь,
					кинематографичный стиль.
	\end{verbatim}
	Субъект: Умный светофор\\
	Среда: перекресток мегаполиса, ночь, дождь\\
	Освещение: (не указано)\\
	Стиль: кинематографичный
	\subsection{журнал промпт-инжиниринга}
	\noindent Промпт v1:
	\begin{verbatim}
					Умный светофор на перекрестке мегаполиса, ночь, дождь,
					кинематографичный стиль.
	\end{verbatim}
	\newpage
	\noindent Результат:
	\begin{figure}[H]
	\centering
	\includegraphics[width=15cm]{z1/1.png}
	\caption{Промпт v1, модель (1)}
	\end{figure}
	\noindent Анализ ошибок: не хватает рамок выделения \\
	\noindent Корректировка: уточнить запрос
	\newpage
	\noindent Промпт v2:
	\begin{verbatim}
					Центральный объект: Умный светофор с камерами.
					Обязательные детали: Объективы камер наблюдения, рамки
					выделения объектов (bounding boxes) вокруг машин.
					Окружение / Фон: Перекресток мегаполиса под дождем.
					Стиль изображения: Кинематографичный кадр (как в кино).
					Освещение: Ночное уличное, отражения на мокром асфальте.
	\end{verbatim}
	\noindent Результат:
	\begin{figure}[H]
	\centering
	\includegraphics[width=15cm]{z1/2.png}
	\caption{Промпт v2, модель (2)}
	\end{figure}
	На изображении видно, что модель задала рамкам разные цвета и подписи рамок (там где они присутсвуют) не всегда соответсвуют изображению. \\
	\newline
	\noindent Анализ ошибок: ненужный дисплей на светофоре\\
	\noindent Корректировка: указать в негативном промпте дисплей на светофоре\\
	\newline
	\noindent Промпт v3:
	\begin{verbatim}
					Центральный объект: Умный светофор с камерами.
					Обязательные детали: Объективы камер наблюдения, рамки
					выделения объектов (bounding boxes) вокруг машин.
					Исключить все остальные детали с Центрального объекта,
					кроме обязательных деталей.
					Окружение / Фон: Перекресток мегаполиса под дождем.
					Стиль изображения: Кинематографичный кадр (как в кино).
					Освещение: Ночное уличное, отражения на мокром асфальте.
	\end{verbatim}
	\newpage
	\noindent Результат:
	\begin{figure}[H]
	\centering
	\includegraphics[width=15cm]{z1/3.png}
	\caption{Промпт v3, модель (2)}
	\end{figure}
	\noindent Анализ ошибок: сцена в целом соответствует запросу\\
	\noindent Корректировка: не требуется\\
	\newpage
	\subsection{Сравнение моделей:}
	\begin{table}[h!]
		\centering
		\begin{tabular}{|l|p{5cm}|p{5cm}|}
			\hline
			\textbf{Критерий} & \textbf{модель (1)} & \textbf{модель (2)} \\
			\hline
			Соответствие промпту & 5 & 5 \\
			\hline
			Отсутствие артефактов & 1 & 4 \\
			\hline
			\textbf{Итог} & 3 & 4.5\\
			\hline
		\end{tabular}
	\end{table}
	Модель (2) лучше справляется с генерацией изображений. Модель (1) часто генерирует артефакты (ломанные линии, 
	кривая перспектива).
	\subsection{Финальный промпт:}
	\begin{verbatim}
					Центральный объект: {CENTRAL}
					Обязательные детали: {REQUIRED}
					Исключить все остальные детали с Центрального объекта,
					кроме обязательных деталей.
					Окружение / Фон: {ENVIRONMENT}
					Стиль изображения: {STYLE}
					Освещение: {LIGHTING}
	\end{verbatim}
	В этом задании промпт писался с помощью "Структурного подхода" + "Постепенное наращивание деталей".

	\stepcounter{section}
	\newpage
	\section*{Задание \thesection : Реверс-инжиниринг промптов (Image-to-Text-to-Image)}
	\subsection{исходные данные}
	\indent \textbf{Референсный объект (Содержание):} Сгоревший серверный блок питания \\
	\indent \textbf{Окружение / Контекст:} Стол в мастерской \\
	\indent \textbf{Стиль референса:} Реалистичное фото последствий аварии \\
	\indent \textbf{Сложность интерпретации (Семантический разрыв):} Передача текстуры копоти и оплавленного пластика \\
	\indent \textbf{Критический визуальный маркер (Обязательно для восстановления):} Обугленные провода и следы гари на металлическом корпусе
	\stepcounter{section}

	\newpage
	\section*{Задание \thesection : Редактирование изображений через In/Out-painting}
	\subsection{исходные данные}
	\indent \textbf{Исходное изображение (Контекст):} Фотография улицы с номерными знаками машин \\
	\indent \textbf{Тип операции:} Inpainting \\
	\indent \textbf{Объект вмешательства (Промпт маски):} Скрытие/Замена номера (Anonymization) \\
	\indent \textbf{Сложность интеграции (Критерий качества):} Естественность размытия или генерация несуществующего номера \\
	\indent \textbf{Сценарий вариативности (Для аугментации):} Разные форматы номеров или полная пикселизация
	\stepcounter{section}

	\newpage
	\section*{Задание \thesection : Создание презентации технического решения с единым визуальным стилем}
	\subsection{исходные данные}
	\indent \textbf{Тема IT-проекта:} Блокчейн-реестр для логистики \\
	\indent \textbf{Целевая аудитория:} Инвесторы в криптовалюты \\
	\indent \textbf{Визуальный стиль (Эстетика):} Киберпанк / Неон (Темная тема) \\
	\indent \textbf{Цветовая палитра:} Черный фон, неоновый фиолетовый и голубой \\
	\indent \textbf{Сквозной визуальный мотив (Якорь стиля):} Светящиеся кубы данных, соединенные линиями
	\subsection{базовый промпт}
	\begin{verbatim}
					Блокчейн реестр для логистики, темная тема, неоновый
					киберпанк, светящиеся кубы данных.
	\end{verbatim}
	Субъект: Блокчейн реестр\\
	Среда: (не указана)\\
	Освещение: (не указано)\\
	Стиль: киберпанк, неоновый
	\stepcounter{section}

	\newpage
	\section*{Задание \thesection : Визуализация технической архитектуры (Text-to-Diagram)}
	\subsection{исходные данные}
	\indent \textbf{Архитектура системы (Тема):} Процесс разработки (DevOps CI/CD) \\
	\indent \textbf{Ключевые компоненты (Узлы схемы):} Код (ветка), Сборка (печь), Тест (микроскоп), Релиз (ракета) \\
	\indent \textbf{Визуальная метафора связей (Потоки):} Железнодорожные пути \\
	\indent \textbf{Стиль визуализации (Композиция):} Линейная схема слева направо \\
	\indent \textbf{Цветовое кодирование зон:} Градиент от фиолетового (Код) к оранжевому (Продакшн)
	\subsection{базовый промпт}
	\begin{verbatim}
					DevOps CI/CD pipeline, code, build, test, release,
					railway tracks, linear layout.
	\end{verbatim}
	Субъект: code, build, test, release\\
	Среда: (не указана)\\
	Освещение: (не указано)\\
	Стиль: railway tracks, linear layout
	\stepcounter{section}

	\newpage
	\section*{Задание \thesection : Генерация синтетического датасета для задачи компьютерного зрения}
	\subsection{исходные данные}
	\indent \textbf{Целевой класс объектов (для детекции):} Спелые яблоки на дереве (АгроТех) \\
	\indent \textbf{Базовое окружение (Инвариант):} Яблоневый сад \\
	\indent \textbf{Вариативные параметры (Что менять в промпте):} Сорт яблок (красные/зеленые), Плотность листвы, Солнечный свет \\
	\indent \textbf{Граничные случаи (Сложные примеры):} Гнилое яблоко, Яблоко в тени листвы, Гроздь яблок
	\subsection{базовый промпт}
	\begin{verbatim}
					Спелые яблоки на дереве, яблоневый сад.
	\end{verbatim}
	Субъект: спелые яблоки\\
	Среда: яблоневый сад\\
	Освещение: (не указано)\\
	Стиль: (не указан)
	\stepcounter{section}

	\newpage
	\section*{Задание \thesection : Структурированное извлечение данных из изображений (Visual ETL)}
	\subsection{исходные данные}
	\indent \textbf{Тип входного изображения:} Фотография ценников на полке магазина \\
	\indent \textbf{Сущности для извлечения (Ключи):} Название продукта, Цена за единицу, Скидка (если есть), Штрих-код (цифры) \\
	\indent \textbf{Сложность / Шум (Проблема для OCR):} Множество объектов в кадре, блики на ценникодержателях \\
	\indent \textbf{Целевая структура JSON (Пример схемы):} \verb|{"products": [{"name": string, "price": float, "discount_percent": int_or_null}]}|
	\stepcounter{section}

	\newpage
	\section*{Заключение}

\end{document}